% BHU PYQ -- Computer Science Introduction to Programming using C (2024-25 NEP)
\documentclass[11pt,a4paper]{article}
\usepackage[a4paper,top=2.5cm,bottom=2.5cm,left=2.8cm,right=2.8cm,headheight=14pt]{geometry}
\usepackage{amsmath,amssymb}
\usepackage[T1]{fontenc}
\usepackage[utf8]{inputenc}
\usepackage{lmodern,microtype,fancyhdr,enumitem,hyperref,lastpage,parskip}

\hypersetup{colorlinks=true,urlcolor=black,linkcolor=black,pdftitle=CSCMJ11 / CSCMN11: Introduction to Programming using C -- BHU NEP 2024-25}

\pagestyle{fancy}\fancyhf{}
\renewcommand{\headrulewidth}{0.4pt}\renewcommand{\footrulewidth}{0.4pt}
\fancyhead[L]{\small\textit{Banaras Hindu University}}
\fancyhead[R]{\small\textit{CSCMJ11 / CSCMN11: Introduction to Programming using C (2024-25)}}
\fancyfoot[C]{\small Page \thepage\ of \pageref{LastPage}}

\newlist{parts}{enumerate}{2}
\setlist[parts,1]{label=(\alph*),leftmargin=*,itemsep=4pt,topsep=3pt}
\setlist[parts,2]{label=(\roman*),leftmargin=*,itemsep=2pt,topsep=2pt}
\newcommand{\pts}[1]{\hfill\textit{[#1]}}

\begin{document}
\begin{center}
  {\large\bfseries BANARAS HINDU UNIVERSITY}\\[2pt]
  {\normalsize U.G. Semester 1 (NEP) Examination, 2024-25}\\[6pt]
  \rule{0.8\linewidth}{0.5pt}\\[6pt]
  {\large\bfseries COMPUTER SCIENCE}\\[4pt]
  {\large\bfseries Paper: CSCMJ11 / CSCMN11 --- Introduction to Programming using C}\\[4pt]
  {\normalsize Time: 2:30 Hours \hfill Full Marks: 70}
\end{center}
\rule{\linewidth}{0.4pt}
\smallskip
\noindent\textit{Instructions: Attempt \textbf{FIVE} questions including Question No.~1, which is compulsory. The figures in the right-hand margin indicate marks.}
\bigskip

\noindent\textbf{Question 1.} Choose the correct options:\pts{$7 \times 2 = 14$}
\begin{parts}
  \item Which of the following is a valid C variable name?
    \begin{parts}
      \item 1variable
      \item variable_1
      \item variable 1
      \item int
    \end{parts}

  \item Which operator is used for dynamic memory allocation in C?
    \begin{parts}
      \item malloc()
      \item alloc()
      \item new
      \item create()
    \end{parts}

  \item What is the size of `int` data type in 32-bit GCC compiler?
    \begin{parts}
      \item 2 bytes
      \item 4 bytes
      \item 8 bytes
      \item 1 byte
    \end{parts}

  \item Which format specifier is used for printing a pointer address?
    \begin{parts}
      \item %d
      \item %f
      \item %p
      \item %s
    \end{parts}
\end{parts}

\bigskip
\noindent\textbf{Question 2.}
\begin{parts}
  \item Explain control structures in C with examples (`if-else`, `switch-case`, loops).\pts{7}
  \item Write a C program to find the factorial of a number using recursion.\pts{7}
\end{parts}

\bigskip
\noindent\textbf{Question 3.}
\begin{parts}
  \item Explain arrays and pointers in C. Write a program to reverse an array using pointers.\pts{7}
  \item Explain structures and unions in C. Differentiate between structure and union with examples.\pts{7}
\end{parts}


\vfill\rule{\linewidth}{0.4pt}
\end{document}
